\documentclass[11pt]{article}

\usepackage[margin=1in]{geometry}
\usepackage{amsmath,amssymb,amsthm}
\usepackage{array,booktabs}
\usepackage{enumitem}
\usepackage[hidelinks]{hyperref}

\theoremstyle{plain}
\newtheorem{observation}{Observation}
\newtheorem{proposition}{Proposition}
\newtheorem{lemma}{Lemma}
\newtheorem{corollary}{Corollary}
\newtheorem{theorem}{Theorem}
\newtheorem*{question}{Open question}
\newtheorem{nogo}{Theorem}
\renewcommand{\thenogo}{\Alph{nogo}}
\theoremstyle{definition}
\newtheorem{definition}{Definition}
\newtheorem{remark}{Remark}

\newcommand{\No}{\mathbf{No}}
\newcommand{\On}{\mathrm{On}_2}
\newcommand{\F}{\mathbb{F}}
\newcommand{\Tr}{\operatorname{Tr}}
\newcommand{\Fr}{\operatorname{Fr}}
\newcommand{\Arf}{\operatorname{Arf}}
\newcommand{\nimmul}{\otimes}
\newcommand{\Cl}{\mathrm{Cl}}
\newcommand{\Sp}{\operatorname{Sp}}
\newcommand{\Orth}{\mathrm{O}}
\newcommand{\AO}{\mathrm{AO}}
\newcommand{\AGL}{\operatorname{AGL}}
\newcommand{\GL}{\operatorname{GL}}
\newcommand{\Aut}{\operatorname{Aut}}
\newcommand{\Inn}{\operatorname{Inn}}
\newcommand{\Out}{\operatorname{Out}}
\newcommand{\rad}{\operatorname{rad}}
\newcommand{\wt}{\operatorname{wt}}
\newcommand{\mex}{\operatorname{mex}}
\newcommand{\Stab}{\operatorname{Stab}}

\title{Gold--Arf Normal Play and the Observation Boundary}
\author{a9lim}
\date{August 2026}

\begin{document}
\maketitle

\begin{abstract}
The Gold trace forms
\[
  Q_a(x)=\Tr(x\,x^{2^a})
\]
are quadratic forms over fields of characteristic $2$. In the nimber setting
they are also composites of combinatorial-game operations:
nim-product via Turning Corners, Frobenius squaring via the diagonal product,
and XOR via disjunctive sum. We construct a uniform finite normal-play game
whose $P$-positions are exactly $\{x:Q_a(x)=0\}$. A public,
refinement-blind Witt frame turns the polar support into a matching. Each
active original coordinate contributes a matched source pair whose overlap
has the corresponding diagonal weight. A FIFO strategy forces zero correction
parity for either seat on every matching-plus-isolates board, and a phase-aware
terminal claim compiles the forced charge to ordinary normal play. The Arf
invariant, together with radical data in the singular case, therefore gives
the exact second-player win bias.

The construction is accompanied by two sharp boundaries. First, any
transcript-stable exact realizer must observe vectors spanning the input; with
weight-$w$ observations, a weight-$t$ input needs at least $t/w$ observations.
Thus a constant total observation budget is impossible, and the construction's
singleton source pairs attain the sharp bound. Second, outcome-preserving
terminal padding gives every finite normal-play tree a refinement-sensitive
winning fork, so fork-based extensional tests cannot certify naturality. The
paper also identifies the quadratic datum with the squaring map of an
extraspecial central extension and proves linear, symmetric, commutative, and
bounded-oracle no-go results that explain why the matching lift is necessary.
\end{abstract}

\section{Introduction and scope}\label{sec:scope}

A short partizan game is built recursively as $\{L\mid R\}$, and games form a
partially ordered abelian group under disjunctive sum~\cite{Conway76,WW}. They
do not form a ring: Conway multiplication is defined on the number subclass,
not on arbitrary games. Since a Clifford algebra needs a commutative scalar
ring, the phrase ``Clifford over games'' has to be interpreted through scalar
subworlds:

\begin{center}
\begin{tabular}{@{}>{\raggedright\arraybackslash}p{0.18\linewidth}
>{\raggedright\arraybackslash}p{0.28\linewidth}
>{\raggedright\arraybackslash}p{0.45\linewidth}@{}}
\toprule
world & algebraic structure & role here \\
\midrule
$\No$ & real-closed field, char $0$ & ordinary real Clifford theory with surreal scalars \\
$\No[i]$ & algebraically closed field, char $0$ & ordinary complex Clifford theory \\
finite nimber subfields & finite fields of char $2$ & characteristic-$2$ Clifford and Arf experiments \\
\bottomrule
\end{tabular}
\end{center}

The implemented \texttt{Nimber} scalar is finite: \texttt{u128} realizes
$\F_{2^{128}}$, containing the nimber subfields $\F_{2^m}$ for
$m=1,2,4,\ldots,128$. It is not the whole proper-class field $\On$.

\section{The char-2 Clifford datum}\label{sec:datum}

In characteristic $2$, a quadratic form is not determined by its polar form.
For basis vectors the Clifford relations are
\[
  e_i^2=q_i,\qquad e_i e_j+e_j e_i=b_{ij}.
\]
The polar form is alternating, so $b_{ii}=0$, but $q_i$ can be nonzero. A
faithful characteristic-$2$ implementation must store $q$ and $b$ separately.
This is the main implementation discipline behind the nimber backend, and --
as Section~\ref{sec:extraspecial} makes precise -- it is also the
group-theoretic content of the whole note: $q$ and $b$ are the independent
central data (squares and commutators) of an extraspecial-type extension.

For a nonsingular quadratic form over $\F_2$ with symplectic basis
$(a_k,b_k)$, the Arf invariant is
\[
  \Arf(Q)=\sum_k Q(a_k)Q(b_k)\in \F_2.
\]
It classifies nonsingular binary quadratic spaces at fixed dimension. Over a
finite characteristic-$2$ field, the corresponding value lies in
$F/\wp(F)$, with $\wp(y)=y^2+y$; for finite fields this quotient is read by the
field trace to $\F_2$.

\begin{remark}
Ovsienko~\cite{Ovsienko16} identifies a precise bridge between Dickson's
classification of quadratic forms over $\F_2$ and the classical classification
of real Clifford algebras. We use that paper as motivation for the Arf/Clifford
connection, not as a blanket classification theorem for all Clifford algebras
over finite nimber fields. In the code, rank and radical data are reported
alongside Arf, and degenerate forms are not collapsed to a single Arf bit.
\end{remark}

The small finite-field check used throughout is:
\[
  q=(\ast2,\ast2),\ b_{01}=\ast1 \quad\text{has Arf }1,
\]
whereas
\[
  q=(\ast2,\ast3),\ b_{01}=\ast1 \quad\text{has Arf }0.
\]
The polar entry is part of the form.

\section{Gold forms from game operations}\label{sec:gold}

Nim-addition is XOR, and XOR is the disjunctive sum of impartial game values.
Nim-multiplication is also game-theoretic: Conway's Turning-Corners product has
the mex recurrence
\begin{equation}\label{eq:tc}
  x\nimmul y
  =
  \mex
  \left\{
    (i\nimmul y)\oplus(x\nimmul j)\oplus(i\nimmul j):
    0\leq i<x,\ 0\leq j<y
  \right\}.
\end{equation}
The repo implements this recurrence directly for small arguments and checks it
against the fast algebraic nim product.

Thus the following operations are game-realizable at the level of nimber
values:
\[
  x\oplus y,\qquad x\nimmul y,\qquad x\mapsto x^2=x\nimmul x,\qquad
  \Tr(x)=x+x^2+\cdots+x^{2^{m-1}}.
\]

For $F=\F_{2^m}$ and $a\geq 1$, define the Gold trace form
\[
  Q_a(x)=\Tr\!\left(x\,x^{2^a}\right)=\Tr\!\left(x^{1+2^a}\right).
\]
Since Frobenius $x\mapsto x^{2^a}$ is additive, $Q_a$ is a quadratic form over
$\F_2$. In the nimber backend, it is also a composite of the game operations
above. Throughout, ``the instance $(m,a)$'' means $Q_a$ on $\F_{2^m}$, and
``$(m,a,\lambda)$'' the scaled component $Q_{\lambda,a}(x)=\Tr(\lambda
x^{1+2^a})$.

\begin{observation}\label{obs:built}
The Gold forms are not merely defined on a field of game values. In the finite
nimber fields tested here, they are built from nim/game operations: Turning
Corners for multiplication, diagonal Turning Corners for Frobenius squaring,
and disjunctive sum for XOR and trace.
\end{observation}

This observation is modest but useful: it gives a concrete quadratic form at
the interface between impartial-game arithmetic and characteristic-$2$
quadratic-form theory.

\subsection{Rank check}

The polar rank of Gold forms is known. In the tested regime
$m=2^k$ and $1\leq a\leq k$, the general formula
\[
  \operatorname{rank} B_{Q_a}=m-\gcd(2a,m)
\]
specializes to $m-2\gcd(a,m)$ because $m/\gcd(a,m)$ is even. The classifier
reproduces the general formula in all tested cases. Representative rows are
shown in Table~\ref{tab:gold}.

\begin{table}[h]
\centering
\begin{tabular}{rrccccc}
\toprule
field & $a$ & $Q$ & $\Arf$ & type & rank & formula \\
\midrule
$\F_{2^{4}}$  & $1$ & $\Tr(x^3)$    & $1$ & $O^-$ & $2$  & $2$ \\
$\F_{2^{8}}$  & $1$ & $\Tr(x^3)$    & $1$ & $O^-$ & $6$  & $6$ \\
$\F_{2^{8}}$  & $2$ & $\Tr(x^5)$    & $1$ & $O^-$ & $4$  & $4$ \\
$\F_{2^{16}}$ & $1$ & $\Tr(x^3)$    & $1$ & $O^-$ & $14$ & $14$ \\
$\F_{2^{16}}$ & $4$ & $\Tr(x^{17})$ & $1$ & $O^-$ & $8$  & $8$ \\
$\F_{2^{32}}$ & $1$ & $\Tr(x^3)$    & $1$ & $O^-$ & $30$ & $30$ \\
\bottomrule
\end{tabular}
\caption{Representative output of \texttt{experiments/trace\_form\_arf.py}.
The script checks all $15$ cases with $m=2,4,8,16,32$ and
$1\leq a\leq \log_2 m$.}
\label{tab:gold}
\end{table}

Most of these forms have a nontrivial radical. The statement is therefore not
``nondegenerate Gold form'' but ``positive-rank quadratic form with a
nonsingular core whose Arf value is computed after symplectic reduction.''

\subsection{A broader trace family}

The coefficient-$1$ Gold form is only one trace monomial. A standard trace
description of the quadratic part of Boolean functions over $\F_{2^m}$ uses
terms
\[
  \Tr(c_i x^{1+2^i})
\]
plus affine terms and, in even degree, a middle half-trace term. The monomials
displayed here are still game-operation-built in the nimber backend:
$c_i x^{1+2^i}=c_i\nimmul x\nimmul x^{2^i}$, followed by trace/XOR.

The current implementation does not claim to package every Boolean quadratic
function. The probe \texttt{experiments/gold\_family\_survey.py} deliberately
omits the middle term and affine bookkeeping. What it does check, exhaustively
over $\F_{256}$, is that scaled Gold components
\[
  Q_{\lambda,a}(x)=\Tr(\lambda x^{1+2^a})
\]
already reach nondegenerate (bent) forms for many $\lambda$. For APN Gold
exponents in that field, the observed count matches the classical
$2(2^m-1)/3$ bent-component count. This matters for the game question
because a bent form has trivial radical: the loopy radical route collapses to
$\{0\}$, and the frame-blind symplectic no-go applies without a degenerate
layer to hide in.

\begin{remark}[a vacuous small target]\label{rem:vacuous}
At $m=4$, $a=1$ the zero set $\{Q_1=0\}\subset\F_{16}$ is the subfield
$\F_4$, an $\F_2$-\emph{subspace} of dimension $2$: the form is genuinely
quadratic (polar rank $2$), but its zero set is affine, so as a target for
``is this $P$-set a genuine quadric'' it is vacuous -- a linear rule realizes
it. The smallest honest unscaled instance is $(8,1)$. The repository's
quadric-fitting layer (\texttt{src/forms/quadric\_fit.rs}) reports exactly
this distinction (genuine polar rank versus affine flat), and every probe
below that consumes a fitted quadric is read with that guard in mind.
\end{remark}

\section{Arf as win-bias}\label{sec:bias}

The relevant counting theorem is classical. If $Q$ is nonsingular on
$\F_2^{2r}$, then
\[
  \#\{x:Q(x)=0\}
  =
  2^{2r-1}+(-1)^{\Arf(Q)}2^{r-1}.
\]
For degenerate forms, the code uses the standard radical-adjusted count: an
anisotropic radical balances the values, while an isotropic radical scales the
bias of the nonsingular core.

The counting theorem becomes the game interpretation after
Theorem~\ref{thm:wittfifo}:

\begin{quote}
For weighted-source \textsc{witt-fifo}, $P$-positions are exactly
$\{x:Q(x)=0\}$. On a nonsingular core the Arf invariant is the sign of the
bias between second-player wins and first-player wins; in the singular case
the radical data supplies the standard adjustment.
\end{quote}

Section~\ref{sec:wittfifo} supplies the rule and proof that make this
interpretation unconditional for the stated game.

\begin{table}[h]
\centering
\begin{tabular}{rrcccr}
\toprule
field & $a$ & $\Arf$ & rank & $\#\{Q=0\}$ & bias \\
\midrule
$\F_{2^{4}}$  & $1$ & $1$ & $2$  & $4$     & $-4$ \\
$\F_{2^{8}}$  & $1$ & $1$ & $6$  & $112$   & $-16$ \\
$\F_{2^{8}}$  & $2$ & $1$ & $4$  & $96$    & $-32$ \\
$\F_{2^{16}}$ & $1$ & $1$ & $14$ & $32512$ & $-256$ \\
$\F_{2^{16}}$ & $4$ & $1$ & $8$  & $30720$ & $-2048$ \\
\bottomrule
\end{tabular}
\caption{Representative zero-counts from
\texttt{experiments/arf\_win\_bias.py}. The bias is
$\#\{Q=0\}-2^{m-1}$.}
\label{tab:bias}
\end{table}

\section{The linear ceiling}\label{sec:linear}

For a normal-play disjunctive sum of impartial games, the P-condition is
linear in Grundy coordinates:
\[
  g_1\oplus\cdots\oplus g_n=0.
\]
So normal-play sums give linear P-sets, not genuine quadrics. The
coin-turning architecture -- the very architecture whose product theorem
defines the nim product -- is subject to the same ceiling.

\begin{nogo}[coin-turning linearity; standard math, machine cross-checked]
\label{thm:A}
Every Winning Ways coin-turning game has configuration $P$-set equal to the
kernel of an $\F_2$-linear nimber-valued map. In particular no coin-turning
$P$-set is a quadric of positive polar rank.
\end{nogo}

\begin{proof}[Proof sketch]
A coin-turning position is a finite set of heads; by the Winning Ways
coin-turning theory~\cite{WW}, the Grundy value of a configuration is the XOR
of the single-coin Grundy values $g(i)$ of its heads. The $P$-set is therefore
$\{x:\bigoplus_{i\in x}g(i)=0\}$, the kernel of the $\F_2$-linear map
$x\mapsto\bigoplus_i x_i\,g(i)$. The repository cross-check
(\texttt{experiments/gold/octal\_attack.py}) verifies XOR-additivity by mex
recursion, exhaustively over all $32{,}768$ companion-family choices on $4$
coins.
\end{proof}

Theorem~\ref{thm:A} has a sharp classical complement: the linear ceiling is
exactly where coding theory already harvests games.

\begin{remark}[the lexicode shadow; standard math and interpretation]
\label{rem:lexicode}
Conway--Sloane lexicodes~\cite{ConwaySloane86} are built by the greedy
lexicographic rule -- the mex rule -- and the 1986 paper realizes them as the
Grundy-value-$0$ positions of associated turning-game move structures; binary
lexicodes are linear \emph{because of} Sprague--Grundy theory (XOR-closure is
a game theorem there, not a coding theorem). They include the Hamming codes
and the length-$24$, $d=8$ extended binary Golay code; over base $2^k$ they
are closed under nim-addition, and they are linear at the Fermat-power bases
$2^{2^k}$ -- exactly the sizes at which nim-multiplication makes the ordinals
below the base a field. Read together with Theorem~\ref{thm:A}: lexicodes are
the \emph{floor} (rich linear codes naturally realized as Grundy-$0$ sets) and
Theorem~\ref{thm:A} is the \emph{ceiling} (coin-turning realizes only linear
$P$-sets); floor and ceiling meet at linear. A quadratic refinement cannot
remain inside this disjunctive-sum ceiling: its obstruction is precisely the
XOR-closure failure measured by the polar form.  Weighted-source FIFO supplies
the required quadratic realization by leaving the disjunctive-sum model.
\end{remark}

For a quadratic form in characteristic $2$,
\[
  Q(u+v)=Q(u)+Q(v)+B(u,v).
\]
Thus for $u,v\in\{Q=0\}$,
\[
  u+v\in\{Q=0\}\quad\Longleftrightarrow\quad B(u,v)=0.
\]
The polar form is exactly the obstruction to XOR-closure. The first two layers
are game-built: XOR is disjunctive sum and $B$ is built from nim products. The
play layer supplied in Section~\ref{sec:wittfifo} integrates that bilinear data
into the quadratic zero set.

\section{Routes and obstructions}\label{sec:blocks}

\subsection{Test benches}\label{sec:benches}

The repo contains several probes. They are useful for narrowing the target, but
they are not a proof that no natural game exists.

\begin{itemize}[leftmargin=*]
\item \texttt{fit\_f2\_quadratic} tests whether a subset of $\F_2^k$ is the zero
set of a quadratic polynomial, and whether the result has nonzero polar rank
(distinguishing genuine quadrics from affine flats; cf.\
Remark~\ref{rem:vacuous}).
\item A bounded misere quotient computer tests small quotient P-sets. The octal
sweep checked $292$ octal codes, heap cutoff at most $4$, and found quotient
orders $2,6,10,12,14$, with no elementary abelian $2$-group of rank at least
$2$. Theorem~\ref{thm:C} below shows this emptiness is forced, not a sampling
accident.
\item \texttt{experiments/misere\_kernel.py} checks the kernel obstruction on
the smallest wild misere quotient $R_8$: the canonical group/kernel shadow is
linear, while the genuinely misere P-element lies outside that vector-space
part. Corollary~\ref{cor:kernel} below shows the linearity is forced by
commutativity.
\item An interactive-kernel solver tests finite move graphs. It confirms that
an ad hoc acyclic game can realize any subset, and that a rule referencing $Q$
itself realizes $\{Q=0\}$ tautologically. The tested $B$-coupled rules do not
produce the Gold zero set.
\item \texttt{examples/loopy\_quadric.rs} tests cyclic move graphs with Draw
sets. The symmetric $B$-coupled rule has Loss-set equal to the radical $R(B)$,
so it reproduces the obstruction rather than the Gold zero set; this is the
special case of Theorem~\ref{thm:G} below, and its small $(4,1)$ ``hit'' is
the radical coincidence $|\{Q_1=0\}|=4=|R(B)|$ there.
\item \texttt{examples/bent\_route.rs} repeats route probes on a bent scaled
Gold component. In that example, a $B$ plus coordinate-frame rule reaches a
quadric of the correct Arf class but not the specific Gold zero set; the naive
local-field/spin-flip completion leaves the quadric variety.
\end{itemize}

Except for the frame-blind proposition below and the lettered theorems of
Section~\ref{sec:nogos}, these are bounded or example-level negative results.

\subsection{Frame-blind no-go}

There is a clean obstruction for rules that see only a nondegenerate
symplectic form and no frame.

\begin{proposition}\label{prop:nogo}
Let $B$ be a nondegenerate alternating form on $V=\F_2^{2r}$ with $r\geq 2$.
If a finite normal-play game has position set $V$ and a move relation invariant
under $\Sp(B)$, then its $P$-set is not a nondegenerate quadric.
\end{proposition}

\begin{proof}
Each element of $\Sp(B)$ is an automorphism of the move graph, so
the retrograde outcome function is $\Sp(B)$-invariant. The P-set
is therefore a union of symplectic orbits. For $r\geq2$, the symplectic group is
transitive on $V\setminus\{0\}$~\cite{Taylor}, so the invariant subsets are only
$\varnothing$, $\{0\}$, $V\setminus\{0\}$, and $V$. A nondegenerate quadric in
dimension at least $4$ has $2^{2r-1}\pm2^{r-1}$ points, which is none of these.
\end{proof}

The hypothesis matters. Coordinate-dependent rules have already broken
$\Sp(B)$ symmetry, so this proposition does not explain every
negative probe. It explains only genuinely frame-blind $B$-only rules.
Section~\ref{sec:nogos} extends the obstruction far beyond this symmetric
class.

\subsection{The diagonal framing}\label{sec:framing}

Once a coordinate frame is allowed, a canonical-looking quadratic refinement of
$B$ is
\[
  Q_{\mathrm{frame}}(v)=\sum_{i<j}B(e_i,e_j)v_i v_j.
\]
For the Gold polar forms tested in this repo, the Gold form decomposes as
\[
  Q_a(v)=Q_{\mathrm{frame}}(v)+\ell_{\mathrm{diag}}(v),
  \qquad
  \ell_{\mathrm{diag}}(v)=\sum_i Q_a(e_i)v_i.
\]
In the tested genuinely quadratic Gold cases up to $\F_{2^{16}}$,
$Q_{\mathrm{frame}}$ has Arf $0$ and the diagonal term changes the form to the
Gold Arf-$1$ class.

\begin{remark}
This is an experimental pattern for the Gold polar forms tested here, not a
theorem about every alternating form plus every coordinate frame. Random
nondegenerate alternating forms can give a zero-diagonal frame quadric of Arf
$1$.
\end{remark}

The diagonal $q_i=Q_a(e_i)$ is therefore the named missing datum. Two results
of the parallel run locate it more precisely.

\begin{proposition}[game-native diagonal source for $a=1$; implemented and
tested]\label{prop:wp}
Let $m=2^k$ and let $w\in\F_{2^m}$ be the XOR of the Fermat coins
$2^{2^t}$ for $t=1,\dots,k-1$. Then the $a=1$ Gold diagonal satisfies
\[
  q_i^{(m,1)}=Q_1(e_i)=\Tr\!\bigl(\wp(w)\,e_i\bigr),
  \qquad
  \wp(w)=w\nimmul w\oplus w,
\]
i.e.\ the trace-dual element of the diagonal is $\wp(w)$. Verified at
$m=4,8,16,32$ via the trace-pairing uniqueness of the dual element
(\texttt{experiments/gold/gold\_diag\_probe.py}).
\end{proposition}

This gives a game-native, tower-uniform, commutative source for the $a=1$
diagonal: one nim-squaring, one XOR, one trace -- all inside the licensed
operation chain of Section~\ref{sec:gold}. The following theorem closes the
same sourcing question for every exponent, including even $a$.

\begin{theorem}[all-exponent diagonal source]\label{thm:all-diagonal-source}
Let $m=2^k\ge2$ and let $\lambda_{a,c}^{(m)}$ be the unique trace-dual of the
scaled diagonal
\[
 \Tr_m(\lambda_{a,c}^{(m)}e_i)
   =\Tr_m(c e_i\Fr^a(e_i)).
\]
For $m=2M$, write $\F_{2^m}=\F_{2^M}(u)$ with
$u=e_M$ and $u^{2^M}=u+1$. Put
\[
 d_0=\Tr_{m/M}(c),\qquad
 d_1=\Tr_{m/M}(c u\Fr^a(u)),\qquad
 L_j=\lambda_{a,d_j}^{(M)}.
\]
Then
\[
 \lambda_{a,c}^{(2M)}=(L_0+L_1)+uL_0.
\]
In particular, for $c=1$ one has $L_0=0$, so
$\lambda_{a,1}^{(2M)}\in\F_{2^M}$ and
$\Tr_m(\lambda_{a,1}^{(2M)})=0$. Consequently there is a
$w\in\F_{2^m}$ with
\[
 \lambda_{a,1}^{(m)}=w^2+w,
 \qquad
 Q_a(e_i)=\Tr_m((w^2+w)e_i)
\]
for every $i$ and every $a$.
\end{theorem}

The proof is the quadratic-tower trace calculation in
\path{writeups/gold_diagonal_source.tex}.  The implementation follows the
same recursion and solves the final Artin--Schreier equation by exact
finite-field linear algebra.  The tower block identities and absolute-trace
descent are kernel-checked in \path{formal/Ogdoad/GoldDiagonal.lean}. Thus
the variation across tower levels is the recursive dual rather than an
obstruction: every unscaled dual has a uniform constructive preimage. The
claim is deliberately unscaled; a general coefficient $c$ need not descend
or have trace zero. The source $w$ is selected by exact finite-field linear
algebra, not by a simple Fermat-coin formula as in Proposition~\ref{prop:wp};
this is a uniform constructive source, not yet a strategic naturality result.

\begin{remark}[subfield vanishing; standard math]
$q_i^{(m,a)}=0$ for all $i<m/2$: the Fermat basis elements at lower indices,
their Frobenius iterates, and their nim products stay inside the half-field,
whose absolute trace from an even-degree extension vanishes. The Gold diagonal
lives on the top Fermat layer only.
\end{remark}

\begin{remark}[diagonal framing]
The diagonal framing
\[
  q_i=Q_a(e_i)=\Tr(e_i^{1+2^a})
\]
is game-built for every $a$: Proposition~\ref{prop:wp} gives the explicit
$a=1$ source, and Theorem~\ref{thm:all-diagonal-source} gives a constructive
source for every exponent. The weighted-source Witt--FIFO theorem in
Section~\ref{sec:wittfifo} supplies a local normal-play rule with
P-set $\{x:Q_a(x)=0\}$.
\end{remark}

\section{The extraspecial reframing}\label{sec:extraspecial}

The diagonal-framing question can be restated in group-theoretic terms, and
the restatement sharpens it. The dictionary is classical -- the
characteristic-$2$ Heisenberg/extraspecial picture of quadratic
forms~\cite{Quillen71,Gorenstein80} -- and this section proves the pieces we
use, to keep the note self-contained. Claim levels: Lemmas~\ref{lem:extdict},
\ref{lem:arftypes}, \ref{lem:abelian} and
Proposition~\ref{prop:between} are standard mathematics; the $R_8$
instantiation in Corollary~\ref{cor:kernel} is implemented and tested
(\path{experiments/misere_kernel.py}); reading $E$-equivariance as a
naturality criterion is interpretation. The weighted-source theorem below
realizes the quadric without claiming that ordinary play itself is a model of
the central extension $E$; that stricter extension-model question is optional.

Throughout, $V=\F_2^n$ written additively, and $Z=\{1,z\}\cong\mathbb{Z}/2$
written multiplicatively. A \emph{quadratic map} $Q\colon V\to\F_2$ is a
function with $Q(0)=0$ whose polarization
$B(x,y)=Q(x+y)+Q(x)+Q(y)$ is bilinear; $B$ is then alternating (hence
symmetric), and $Q$ may have any diagonal $q_i=Q(e_i)$, exactly the
char-$2$ datum the engine keeps separate from $b$.

\subsection{Quadratic forms are central extensions}

\begin{lemma}[extraspecial dictionary]\label{lem:extdict}
Let $1\to Z\to E\xrightarrow{\ \pi\ }V\to 1$ be a central extension of the
elementary abelian $2$-group $V$ by $Z\cong\mathbb{Z}/2$.
\begin{enumerate}[label=(\roman*),leftmargin=*]
\item For $x\in V$ and any lift $\tilde x\in\pi^{-1}(x)$ one has
$\tilde x^2\in Z$, and $\tilde x^2$ is independent of the lift; the
\emph{squaring form} $Q_E\colon V\to\F_2$ defined by
$\tilde x^2=z^{Q_E(x)}$ is well defined, as is the \emph{commutator pairing}
$B_E$ defined by $[\tilde x,\tilde y]=z^{B_E(x,y)}$, which is bilinear and
alternating. They satisfy
\[
(\tilde x\tilde y)^2=\tilde x^2\,\tilde y^2\,[\tilde x,\tilde y],
\qquad\text{equivalently}\qquad
Q_E(x+y)=Q_E(x)+Q_E(y)+B_E(x,y),
\]
so $Q_E$ is a quadratic map with polar form $B_E$.
\item Conversely, every quadratic map $Q$ arises: fixing a basis
$e_1,\dots,e_n$ of $V$, the bilinear $2$-cocycle
\[
\varphi(x,y)=\sum_i x_iy_i\,Q(e_i)+\sum_{i>j}x_iy_j\,B(e_i,e_j)
\]
defines a group $E_Q=Z\times V$ with
$(s,x)(t,y)=(s+t+\varphi(x,y),\,x+y)$, whose squaring form is $Q$ and whose
commutator pairing is the polar form $B$.
\item The assignment $Q\mapsto[E_Q]$ is a bijection from quadratic maps on
$V$ to equivalence classes of central extensions of $V$ by $\mathbb{Z}/2$.
\item $E_Q$ is abelian iff $B=0$ (iff $Q$ is linear); and for $n\geq 1$,
$E_Q$ is \emph{extraspecial} (i.e.\ $Z(E)=[E,E]=\Phi(E)\cong\mathbb{Z}/2$)
iff $B$ is nondegenerate, which forces $n=2r$ even and $|E_Q|=2^{1+2r}$.
\end{enumerate}
\end{lemma}

\begin{proof}
(i) $\pi(\tilde x^2)=2x=0$, so $\tilde x^2\in Z$; replacing $\tilde x$ by
$z\tilde x$ multiplies the square by $z^2=1$. Commutators of lifts lie in
$Z$ because $V$ is abelian, are central, and are unchanged by central
retagging of the lifts; centrality of commutators gives bimultiplicativity,
and $[\tilde x,\tilde x]=1$ gives $B_E(x,x)=0$. For the identity: with
$c=[\tilde y,\tilde x]$ central,
$\tilde x\tilde y\tilde x\tilde y
=\tilde x(\tilde x\tilde y\,c)\tilde y=\tilde x^2\tilde y^2c$,
and $c=c^{-1}=[\tilde x,\tilde y]$ since $Z$ has exponent $2$. As
$\tilde x\tilde y$ lifts $x+y$, the additive identity follows.

(ii) A bilinear map is a $2$-cocycle, so $E_Q$ is a group with center
containing $Z=\{(0,0),(1,0)\}$. Squares: $(s,x)^2=(\varphi(x,x),0)$ and
$\varphi(x,x)=\sum_i x_iQ(e_i)+\sum_{i>j}x_ix_jB(e_i,e_j)=Q(x)$ by the
polarization expansion of $Q$ in the basis. Commutators:
$(s,x)(t,y)(s,x)^{-1}(t,y)^{-1}=(\varphi(x,y)+\varphi(y,x),0)
=(B(x,y),0)$, using that the symmetrization of $\varphi$ is $B$ (the
diagonal terms cancel and $B$ is symmetric).

(iii) Equivalent extensions have cohomologous cocycles, and a coboundary
$\delta\psi(x,y)=\psi(x)+\psi(y)+\psi(x+y)$ has zero diagonal
($\delta\psi(x,x)=\psi(0)=0$), so the squaring form is an invariant of the
class; (ii) gives surjectivity. Injectivity: if two extensions have equal
squaring form $Q$, choose sections and cocycles $\varphi,\varphi'$; both
have diagonal $Q$, and both have symmetrization equal to the polar form of
$Q$, so $d=\varphi+\varphi'$ is a symmetric cocycle with zero diagonal. The
extension $E_d$ is then abelian of exponent $2$, i.e.\ an $\F_2$-vector
space, so $1\to Z\to E_d\to V\to 1$ splits and $d$ is a coboundary; the two
extensions are equivalent.

(iv) Commutators generate $z^{\,\mathrm{im}\,B}$, so $E_Q$ is abelian iff
$B=0$. In general $Z(E_Q)=\pi^{-1}(\rad B)$, since $\tilde x$
is central iff $B(x,\cdot)=0$; thus $Z(E_Q)=Z$ iff $B$ is nondegenerate. If
so, $B\neq0$ gives $[E,E]=Z$, and $\Phi(E)=E^2[E,E]=Z$ because all squares
lie in $Z$ and $E/Z$ is elementary abelian. A nondegenerate alternating
form exists only in even dimension.
\end{proof}

\begin{remark}
This is the characteristic-$2$ Heisenberg picture: $E_Q$ is the Heisenberg
group of $(V,B)$ and $Q$ selects the central extension; cohomologically,
(iii) is the standard identification of $H^2(V;\F_2)$ with quadratic maps
$V\to\F_2$~\cite{Quillen71}. The same finite-quadratic-module data drives
the Weil $S$/$T$ matrices in the library's integral layer. Note that the
dictionary is exactly the engine's discipline in group clothing: the
diagonal $q_i$ (squares) and the off-diagonal $b_{ij}$ (commutators) are
independent central data, and collapsing them is collapsing $E_Q$ onto an
abelian quotient.
\end{remark}

\subsection{Arf classifies the two extraspecial types}

Write $H$ for the hyperbolic plane ($Q(e_1)=Q(e_2)=0$, $B(e_1,e_2)=1$;
$\Arf=0$) and $A$ for the anisotropic plane ($Q\equiv1$ on $V\setminus\{0\}$;
$\Arf=1$), and $\circ$ for the central product (identify the centers).

\begin{lemma}[Arf = the two types]\label{lem:arftypes}
Let $B$ be nondegenerate, $\dim V=2r\geq2$.
\begin{enumerate}[label=(\roman*),leftmargin=*]
\item $E_{Q\perp Q'}\cong E_Q\circ E_{Q'}$.
\item $E_H\cong D_8$ and $E_A\cong Q_8$.
\item Every extraspecial $2$-group of order $2^{1+2r}$ is isomorphic to
$E_Q$ for a nondegenerate $Q$, and
$E_Q\cong E_{Q'}$ iff $(V,Q)\cong(V',Q')$ iff the ranks and Arf invariants
agree~\cite{Dickson01,Arf41}. Hence there are exactly two extraspecial
groups of order $2^{1+2r}$:
\begin{align*}
E^+_{2^{1+2r}}&=D_8^{\circ r}
  &&(\Arf=0,\ Q\cong H^{\perp r}),\\
E^-_{2^{1+2r}}&=D_8^{\circ(r-1)}\circ Q_8
  &&(\Arf=1,\ Q\cong H^{\perp(r-1)}\perp A).
\end{align*}
In particular $D_8\circ D_8\cong Q_8\circ Q_8$, the group avatar of
$H\perp H\cong A\perp A$ (additivity of $\Arf$).
\item (Census.) With $\varepsilon=(-1)^{\Arf Q}$,
\[
\#\{g\in E_Q: g^2=1\}=2\,\#\{x:Q(x)=0\}=2^{2r}+\varepsilon\,2^{r},
\]
so the element-order census of $E_Q$ is the zero-count bias of
Section~\ref{sec:bias} read multiplicatively, and it distinguishes the two
types.
\end{enumerate}
\end{lemma}

\begin{proof}
(i) The cocycle $\varphi\oplus\varphi'$ on $V\oplus V'$ is bilinear with
diagonal $Q\perp Q'$, and
$(s,(x,x'))\mapsto[((s,x),(0,x'))]$ is an isomorphism onto
$(E_Q\times E_{Q'})/\langle(z,z')\rangle$.

(ii) $E_H$ has order $8$, is nonabelian ($B\neq0$), and has a noncentral
involution (any lift of $e_1$ squares to $z^{Q(e_1)}=1$), so $E_H\cong D_8$;
in $E_A$ every noncentral element squares to $z$ (as $Q\equiv1$ off $0$), so
$z$ is the unique involution and $E_A\cong Q_8$.

(iii) For $E$ extraspecial, $V=E/Z(E)$ is elementary abelian (as
$\Phi(E)=Z(E)$), so Lemma~\ref{lem:extdict} applies to
$1\to Z(E)\to E\to V\to1$ and gives $E\cong E_{Q_E}$ with $B_{Q_E}$
nondegenerate ($Z(E)$ exactly central). If $\psi\colon E_Q\to E_{Q'}$ is an
isomorphism, it preserves the (characteristic) centers, fixes $z\mapsto z'$
(the unique nontrivial central element), and induces an $\F_2$-linear map
$\bar\psi\colon V\to V'$ with
$Q'(\bar\psi x)$ read from $\psi(\tilde x)^2=\psi(\tilde x^2)$ -- an
isometry. Conversely an isometry $g$ transports the cocycle:
$\varphi'\circ(g\times g)$ has diagonal $Q'\circ g=Q$, hence by
Lemma~\ref{lem:extdict}(iii) yields an extension equivalent to $E_Q$, and
$(s,x)\mapsto(s,gx)$ closes the isomorphism. Dickson's classification of
nonsingular quadratic forms over $\F_2$ (two isometry classes per even
rank, separated by $\Arf$) finishes; the central-product normal forms
follow from (i), (ii) and Witt decomposition.

(iv) The two lifts of $x$ both square to $z^{Q(x)}$, so they are
involutions or the identity iff $Q(x)=0$; now apply the zero-count formula
of Section~\ref{sec:bias}. The counts differ for $\varepsilon=\pm1$, so
$E^+\not\cong E^-$.
\end{proof}

\subsection{The abelian obstruction}

\begin{lemma}[abelian obstruction]\label{lem:abelian}
Let $M$ be a commutative monoid, $Z=\{1,z\}\subseteq M$ a two-element
subgroup, $N\subseteq M$ a submonoid containing $Z$, and
$\pi\colon N\twoheadrightarrow V$ a surjective monoid homomorphism onto an
$\F_2$-vector space with $\pi^{-1}(\pi(m))=mZ$ for all $m\in N$. Then the
squaring form $Q\colon V\to\F_2$, defined by $\tilde m^2=z^{Q(x)}$ for any
$\tilde m\in\pi^{-1}(x)$, is well defined and \emph{$\F_2$-linear}; its
polar form vanishes identically. Consequently no commutative monoid
realizes, through its own squaring map, a quadratic form with nonzero --
a fortiori nondegenerate -- characteristic-$2$ polar form. Equivalently, by
Lemma~\ref{lem:extdict}(iv): if $B\neq0$, the extension $E_Q$ is
nonabelian and admits no model $(N,Z,\pi)$ inside any commutative monoid.
\end{lemma}

\begin{proof}
$\pi(m^2)=2\pi(m)=0$, so $m^2\in\pi^{-1}(0)=Z$, and
$(zm)^2=z^2m^2=m^2$, so $Q$ is well defined. Commutativity gives
$(mn)^2=m^2n^2$; since $\tilde m\tilde n$ is a preimage of $x+y$, this is
$Q(x+y)=Q(x)+Q(y)$, i.e.\ $B\equiv0$.
\end{proof}

\begin{corollary}[misère quotients: the kernel shadow is linear]\label{cor:kernel}
Let $\mathcal{Q}$ be a finite misère quotient -- a finite commutative
monoid -- with kernel $K$, the maximal subgroup of $\mathcal{Q}$ (the
mutual-divisibility class of the product of all idempotents)~\cite{PS}.
\begin{enumerate}[label=(\roman*),leftmargin=*]
\item Every configuration $(N,Z,\pi)$ as in Lemma~\ref{lem:abelian} inside
$\mathcal{Q}$ has linear squaring form: the intrinsic multiplication of a
misère quotient cannot supply a quadratic refinement with $B\neq0$, on $K$
or anywhere else in $\mathcal{Q}$.
\item Under the regularity hypothesis of \cite[Thm.~6.4]{PS} (satisfied by
the regular finite quotients arising in practice), $K\cong(\mathbb{Z}/2)^k$
is the group of normal-play Grundy values and the $P$-portion meets $K$ in
the XOR-linear normal-play set.
\item On the smallest wild quotient
$R_8=\langle a,b,c\mid a^2=1,\ b^3=b,\ bc=ab,\ c^2=b^2\rangle$ with
$P=\{a,b^2\}$: the idempotents are $\{1,b^2\}$, the kernel is
$K=\{b^2,b,ab,ab^2\}\cong(\mathbb{Z}/2)^2$ with identity $b^2$,
$P\cap K=\{b^2\}\mapsto\{0\}$ is linear, and the genuinely misère
$P$-element $a$ lies outside $K$. (Implemented and tested:
\path{experiments/misere_kernel.py}.)
\end{enumerate}
So the linear obstruction observed on $R_8$ is forced by commutativity, not
an accident of the example: a nondegenerate polar form is the commutator
pairing of a nonabelian group, and a misère quotient has none to offer.
\end{corollary}

\begin{proof}
(i) is Lemma~\ref{lem:abelian} applied to submonoids of $\mathcal{Q}$,
which are commutative. (ii) is quoted from \cite{PS}. (iii) is a finite
verification.
\end{proof}

\subsection{The integral shadow}\label{sec:integral}

The same extension data has an integral-lattice avatar, already shipped in
the library (interpretation, verified through the
\texttt{DiscriminantForm} + \texttt{arf\_nimber} pipeline): the Gold
nonsingular core of rank $2r$ and Arf $\varepsilon$ is the discriminant form
of the even lattice $U(2)^{r}$ ($\varepsilon=0$) or
$U(2)^{r-1}\oplus D_4$ ($\varepsilon=1$); the Milgram phase is
$4\cdot\Arf\bmod 8$; the Weil $S$-matrix prefactor is $(-1)^{\Arf}$; and the
Frobenius--Schur indicator of the Heisenberg representation of $E_Q$ is
$(-1)^{\Arf}$. The metaplectic relation $(ST)^3=S^2$ holds iff the diagonal
$T$ uses a quadratic refinement of the correct Arf class -- it selects the
Arf \emph{class}, not the member. So the integral bridge carries exactly one
game-relevant bit; it cannot, by itself, distinguish the Gold form from any
other refinement in its class.

\subsection{\texorpdfstring{An $E$-equivariance screen}{An E-equivariance screen}}

Proposition~\ref{prop:nogo} kills rules that are blind to everything but
the symplectic structure (Tier~1), while per-position evaluator circuits
for $Q_a$ realize the quadric tautologically (Tier~3). The extraspecial
dictionary suggests a criterion strictly between the two: the rule should
see the extension $E_Q$ -- which carries the diagonal data $q_i$
structurally, as squares -- but only up to its automorphisms, so that no
basis, field structure, or evaluation circuit can be smuggled in.

\begin{definition}[uniform rules; the three tiers]\label{def:tier2}
Let $Q$ be nondegenerate on $V=\F_2^{2r}$ with polar form $B$, and
$E=E_Q$, $\pi\colon E\to V$, $\Sigma=\{x\in V: Q(x)=0\}$.
\begin{enumerate}[label=(\alph*),leftmargin=*]
\item A \emph{uniform rule on a finite set $X$} is a single move relation
$M\subseteq X\times X$, fixed once for all positions; for normal play we
require the move digraph acyclic, and outcomes are computed by the usual
retrograde recursion. (For loopy or misère readings, replace the outcome
recursion; the equivariance constraint below applies verbatim, since
outcome classes are invariants of digraph automorphisms.)
\item A uniform rule on $E$ is \emph{$E$-equivariant} if every
$\alpha\in\Aut(E)$ is an automorphism of the move digraph.
Its \emph{shadow} is $\pi(P)\subseteq V$, where $P$ is its $P$-set; it
\emph{realizes} $T\subseteq V$ if $\pi(P)=T$.
\item \emph{Tier 1} ($\Sp(B)$-blind): a uniform rule on $V$
invariant under $\Sp(B)$, as in
Proposition~\ref{prop:nogo}. \emph{Tier 3} ($Q_a$-evaluation): a family of
games $\{\Gamma_x\}_{x\in V}$ whose structure varies with the designated
input $x$ (e.g.\ the trace-circuit evaluator); this is not a uniform rule.
\emph{Tier 2 screen}: a route to the Gold quadric is Tier-2 natural
\emph{only if} its positions and moves lift to an $E$-equivariant uniform
rule on $E_{Q_a}$ -- with $Q_a$ restricted to its nonsingular core when the
form is degenerate, matching the classifier's radical discipline -- that
realizes $\{Q_a=0\}$.
\end{enumerate}
\end{definition}

\begin{proposition}[the screen sits strictly between the solved tiers]\label{prop:between}
Let $Q$ be nondegenerate on $V=\F_2^{2r}$, $r\geq2$, and $E=E_Q$.
\begin{enumerate}[label=(\roman*),leftmargin=*]
\item The image of $\Aut(E)\to \GL(V)$ is exactly the
orthogonal group $\Orth(Q)$, with kernel the central twists
$g\mapsto g\,z^{\ell(\pi g)}$, $\ell\in V^{*}$, which equal
$\Inn(E)\cong V$ because $B$ is nondegenerate; thus
$\Out(E)\cong \Orth(Q)$ (cf.~\cite{Winter72}). The
$\Aut(E)$-orbits on $E$ are
\[
\{1\},\qquad \{z\},\qquad \pi^{-1}(\Sigma\setminus\{0\}),\qquad
\pi^{-1}(V\setminus\Sigma).
\]
\item Hence the $P$-set of any $E$-equivariant uniform rule is a union of
these four orbits, and its shadow is one of the eight unions of $\{0\}$,
$\Sigma\setminus\{0\}$, $V\setminus\Sigma$. The quadric $\Sigma$ is among
them: the Tier-1 exclusion does not apply.
\item (Tier 1 $\subsetneq$ Tier 2.) Every $\Sp(B)$-invariant
uniform rule on $V$ pulls back along $\pi$ to an $E$-equivariant rule with
$P$-set $\pi^{-1}$ of the original, whose shadow is therefore one of
$\varnothing$, $\{0\}$, $V\setminus\{0\}$, $V$ -- never the quadric
(Proposition~\ref{prop:nogo} in this language). The containment is strict:
the \emph{squaring rule}
\[
g\to h\ \text{legal}\iff g^2=z\ \text{and}\ h^2=1
\]
(``move from any order-$4$ position to any position of order at most
$2$'') is $E$-equivariant and acyclic, and its $P$-set is the involution
locus $\{g:g^2=1\}=\pi^{-1}(\Sigma)$, with shadow exactly $\Sigma$.
\item (Tier 2 $\subsetneq$ Tier 3.) Up to the isometry induced by any group
isomorphism, the family of $E$-equivariant rules and their shadows depends
only on $(r,\Arf Q)$ (Lemma~\ref{lem:arftypes}(iii)). In particular the
screen cannot separate Gold forms of equal core rank and Arf invariant, and
an $E$-equivariant rule has no access to $m$, the exponent $a$, the field
multiplication, the coordinate frame $q_i=Q_a(e_i)$, or any evaluation
circuit. Conversely, arbitrary subsets of $V$ -- all realizable by ad hoc
acyclic lookup games and by Tier-3 evaluators
(Section~\ref{sec:benches}) -- are not shadows of $E$-equivariant rules once
they leave the eight-set list of (ii).
\end{enumerate}
\end{proposition}

\begin{proof}
(i) Automorphisms preserve $Z=Z(E)$ and fix $z$
($\Aut(\mathbb{Z}/2)=1$), so the induced map preserves the
squaring form: $Q(\bar\alpha x)$ is read from
$\alpha(\tilde x)^2=\alpha(\tilde x^2)=z^{Q(x)}$; the image lies in
$\Orth(Q)$. Surjectivity: an isometry $g$ transports the standard cocycle as in
the proof of Lemma~\ref{lem:arftypes}(iii), producing an automorphism of
$E_Q$ inducing $g$. The kernel consists of maps $g\mapsto g\,z^{\ell(\pi g)}$
with $\ell$ linear (multiplicativity forces additivity of $\ell$); inner
automorphisms are exactly the twists $\ell=B(\bar g,\cdot)$, and
nondegeneracy makes $\bar g\mapsto B(\bar g,\cdot)$ onto $V^{*}$. Orbits:
$\Orth(Q)$ is transitive on $\Sigma\setminus\{0\}$ and on $V\setminus\Sigma$ by
Witt's extension theorem in its characteristic-$2$ quadratic-space form
\cite{Taylor} (both sets are nonempty for $r\geq2$); the two lifts of any
$x\neq0$ are conjugate under conjugation by $\tilde g$ with $B(\bar g,x)=1$;
and $1$, $z$ are fixed by everything.

(ii) Outcome recursion is invariant under digraph automorphisms.  Thus $P$
is $\Aut(E)$-invariant; apply (i) and project.

(iii) The pullback $M=\{(g,h):(\pi g,\pi h)\in R\}$ of an
$\Sp(B)$-invariant rule $R$ is preserved by
$\Aut(E)$, since the induced action on $V$ lies in
$\Orth(Q)\subseteq\Sp(B)$; it is acyclic when $R$ is, and an easy
induction along the acyclic rank gives
$P(M)=\pi^{-1}(P(R))$. By the transitivity of
$\Sp(B)$ on $V\setminus\{0\}$ for $r\geq2$, $P(R)$ is a union
of $\{0\}$ and $V\setminus\{0\}$. For the squaring rule: automorphisms
preserve element orders, so the rule is $E$-equivariant; positions with
$g^2=1$ are terminal, hence $P$; positions with $g^2=z$ have a move (to
$1$), hence $N$. Its shadow $\Sigma$ satisfies
$1<|\Sigma|=2^{2r-1}+\varepsilon2^{r-1}<2^{2r}-1$ for $r\geq2$, so it is
none of the four Tier-1 shadows.

(iv) If $\psi\colon E_Q\to E_{Q'}$ is an isomorphism (it exists iff ranks
and Arf agree, Lemma~\ref{lem:arftypes}(iii)), then
$M\mapsto(\psi\times\psi)(M)$ is a bijection between $E$-equivariant rules
carrying $P$-sets to $P$-sets and shadows to shadows through the induced
isometry $\bar\psi$; in particular it carries rules realizing $\{Q=0\}$ to
rules realizing $\{Q'=0\}$. The shadow constraint of (ii) excludes all but
eight subsets, while Tier-3 constructions realize any subset; the
containment is proper.
\end{proof}

\begin{remark}[what the screen does and does not settle]\label{rem:screen}
Equivariance is a screen, not a construction. Over the \emph{abstract}
group $E$ the screen is satisfiable -- Proposition~\ref{prop:between}(iii)
exhibits the squaring rule, which is nothing but the extraspecial squaring
map of Lemma~\ref{lem:extdict} read as a one-move relation. That rule
consults only the multiplication of $E$, so it is a $Q$-evaluator exactly
insofar as $E$ itself is fed in by hand. The reframing replaces the request
\emph{find a play rule that computes $Q_a$} by the structural question
\emph{exhibit a game-native model of the extension $E_{Q_a}$} --
positions built from game values, multiplication from game constructions --
after which the rule layer is canonical. Lemma~\ref{lem:abelian} is the
sharp constraint on that search: no commutative value-world, in particular
no misère-quotient kernel (Corollary~\ref{cor:kernel}), can host $E$ once
$B\neq0$. The noncommutativity must enter from a structurally available
asymmetry, and the one that normal, misère, and partizan play all possess
-- and that the symmetric polar form $B$ discards -- is the
first-/second-player asymmetry of the move relation. Claim
levels: the proposition is standard mathematics; treating
$E$-equivariance as \emph{the} naturality criterion is interpretation. The
weighted-source rule resolves the P-set problem covariantly without first
constructing the entire abstract extension as a game object; such a full model
remains a stricter optional generalization. (For $r=1$ the screen degenerates:
on the anisotropic plane $\Sigma=\{0\}$ and $\Orth(Q)=\Sp(B)$,
so the statement is vacuous there; the Gold targets of interest have
$r\geq2$ cores.) Theorem~\ref{thm:E} below sharpens the screen further:
read as \emph{invariance} it admits no live middle at all, so the
$E$-structure must enter \emph{covariantly}, in the dynamics.
\end{remark}

\begin{question}
Does any game-native construction realize the extraspecial extension
$E_{Q_a}$ -- equivalently, produce the diagonal data $q_i=Q_a(e_i)$ as
squares in a noncommutative game-built structure -- without an evaluation
circuit for $Q_a$? By Lemma~\ref{lem:abelian} the source cannot be any
commutative game-value monoid.
\end{question}

\section{No-go theorems for the middle tier}\label{sec:nogos}

The results below delimit commutative provenance, affine and orthogonal
symmetry, bounded $B$-oracle access, $B$-local flip rules, and
$E$-translation dynamics. Theorem~\ref{thm:A} supplies the preceding
coin-turning linearity result.

\begin{nogo}[commutative squaring collapse; standard math]\label{thm:B}
In any commutative monoid, squaring is an endomorphism, so any quadratic
datum read through homomorphic coordinates has polarization identically
zero. In particular the extraspecial squaring map $Q$ (for $B\neq0$) cannot
live in, map into, or pull back through any commutative game monoid. This
covers misère quotients, turn-parity counters, and all coin-turning
configuration sums.
\end{nogo}

\begin{proof}
Immediate from Lemma~\ref{lem:abelian}: commutativity gives
$(mn)^2=m^2n^2$, so the squaring form is additive and $B\equiv0$.
\end{proof}

\begin{nogo}[group misère quotients are tiny; conditional]\label{thm:C}
Assume \cite[Thm.~6.4]{PS} with its regularity hypothesis. If the misère
indistinguishability quotient $\mathcal{Q}(\mathcal{A})$ of a set of
impartial games is a group, then $|\mathcal{Q}(\mathcal{A})|\leq 2$.
\end{nogo}

\begin{proof}[Proof sketch]
Doubling lemma: in a group quotient every element is cancellable, and any
position of Grundy value $2$ doubles to a misère-$P$ position,
contradicting Grundy-distinctness from $0$; the quoted structure theorem
then pins the quotient to at most $\{1,a\}$.
\end{proof}

\begin{remark}[scope and a probe caveat]\label{rem:thmC}
Theorem~\ref{thm:C} is conditional on the Plambeck--Siegel structure
theorem as cited; the regularity hypothesis is quoted from the repository's
\path{experiments/misere_kernel.py} docstring, and the JCTA 2008 paper has
\emph{not} been independently re-verified within this program (closing that
gate is a listed next move). Two consequences for the probe map: the
$(\mathbb{Z}/2)^k$ ($k\geq2$) target of \texttt{examples/octal\_hunt.rs} is
empty a priori -- its empirical negatives (quotient orders
$2,6,10,12,14$) are a theorem, not a sampling accident -- and any future
``hit'' from that sweep would be a labeling artifact, since the helper
\texttt{p\_set\_as\_f2} checks only that the subset labeling is a bijection,
not that it is a monoid homomorphism. The hunt's target should be reframed
accordingly.
\end{remark}

For the remaining results, fix a nondegenerate quadratic form $Q$ on
$V=\F_2^{2r}$, $r\geq2$, with polar form $B$ and zero set
$\Sigma=\{Q=0\}$, and write
$\AO(Q)=\{x\mapsto gx+c: g\in\Sp(B),\ Q(gx+c)=Q(x)\ \forall x\}$ for the
affine stabilizer appearing below.

\begin{nogo}[affine symmetry budget]\label{thm:D}
$\Stab_{\AGL(V)}(\Sigma)=\AO(Q)$; it contains no nontrivial pure
translation (for a nonsingular core, $c\in\rad(B)=0$), and its pure-linear
part is exactly $\Orth(Q)$. Consequently every uniform rule whose move digraph
admits even one affine automorphism outside $\AO(Q)$ has $P$-set
$\neq\Sigma$.
\end{nogo}

\noindent\emph{Verification.} An exhaustive check covers all $322{,}560$
affine maps of $\F_2^4$ at $m=4$ in both Arf classes.  The linear-part
identification follows from Witt extension.

\begin{nogo}[the symmetry dial has no middle]\label{thm:E}
$\Orth(Q)$ is maximal in $\Sp(B)$, and the $\Orth(Q)$-orbitals on $V\times V$
are exactly the Gram classes $(Q(u),Q(w),B(u,w))$ together with the
degeneracy flags (Witt extension; see \cite[\S8]{EKM08}). Hence a rule
invariant under any group strictly larger than $\Orth(Q)$ is dead by
Theorem~\ref{thm:D} (its symmetry reaches outside $\AO(Q)$), while a rule
fully $\Orth(Q)$- or $\Aut(E)$-equivariant factors extensionally through
pointwise $Q$- and $B$-evaluation. There is no invariance class strictly
between dead (too coarse) and Gram-factoring (an evaluator): $E$-naturality
must be \emph{covariance}, not invariance.
\end{nogo}

\noindent\emph{Verification.} Maximality verified at $m=4$: every one of
the $648$ (resp.\ $600$) elements of $\Sp(4,2)$ outside $\Orth^{+}(4,2)$
(resp.\ $\Orth^{-}(4,2)$) generates $\Sp(4,2)$ together with the orthogonal
group; orbitals enumerated exhaustively at $m=4$
(\texttt{experiments/gold/ao\_orbitals.py}).

\begin{nogo}[$B$-oracle lower bound]\label{thm:F}
Access model: legality of a move is decided from XOR combinations of the
two positions and $t$ fixed constant vectors, with oracle access to $B$
only. If $t\leq 2r-3$, then for \emph{every} quadratic refinement $Q'$ of
$B$ simultaneously there is a nonzero $Q'$-singular vector in the
complement of the constant span; the transvection it induces is a digraph
automorphism lying outside $\Orth(Q')$, and Theorem~\ref{thm:D} kills the
$P$-set. The bound is tight: at $t=2r-2$ with anisotropic complement the
argument fails. Hence the Gold diagonal framing $q_i=Q_a(e_i)$, which
supplies $2r$ frame constants, is within two dimensions of
information-theoretically forced.
\end{nogo}

\noindent\emph{Verification.} Transvection step and the exact escape
boundary $t=2r-2$ machine-checked at $m=4$; end-to-end spot check on $50$
random $t=1$ $B$-oracle rules, retrograde-solved
(\path{experiments/gold/skeptic_nogo_check.py},
\path{nogo_verify.py}).

\begin{nogo}[$B$-local flip rules]\label{thm:G}
Any rule of the form ``flip $d$ at $v$ iff $f(d,B(v,d))$'' is automatically
undirected, because $B$ is alternating:
$f(d,B(v+d,d))=f(d,B(v,d)+B(d,d))=f(d,B(v,d))$. For undirected loopy games
the outcome partition collapses: $\mathrm{Win}=\varnothing$,
$\mathrm{Loss}=$ the isolated vertices $=$ an affine flat, and
$\mathrm{Loss}\cup\mathrm{Draw}\neq\Sigma$ for $r\geq2$ (the count
$|\Sigma|=2^{r-1}(2^{r}+\varepsilon)$ has an odd factor $>1$). This covers
the Loss, Draw, and $\mathrm{Loss}\cup\mathrm{Draw}$ targets at once. The
symmetric-$B$ rule of \texttt{examples/loopy\_quadric.rs} is the special
case $f(d,\beta)=\beta$; its $(4,1)$ ``hit'' is the radical coincidence of
Section~\ref{sec:benches}.
\end{nogo}

\noindent\emph{Verification.} Undirectedness is the two-line identity
above (standard math); the loopy outcome collapse including the Draw-set
gap is machine-checked at $m=4$
(\path{experiments/gold/skeptic_nogo_check.py}).

\begin{nogo}[$E$-translation no-go]\label{thm:H}
For $r\geq2$, no left-$E$-equivariant rule on $E=E_Q$ with moves
$g\to gn$, $n\in N$ for a fixed nonempty $N\subseteq E$, has $P$-set equal
to the involution locus $I=\{g:g^2=1\}=\pi^{-1}(\Sigma)$.
\end{nogo}

\begin{proof}
$I\cdot I=E$: for every $v\neq0$, the count
$\#\{x:Q(x)=0=Q(x+v)\}=2^{2r-2}+(-1)^{\Arf Q}2^{r-1}$ is positive for
$r\geq2$ (character sum), so every $n\in N$ factors as $n=gh$ with
$g,h\in I$. Then $g\in I$, and $g\cdot n=g(gh)=g^2h=h\in I$: the move
$g\to gn$ is an edge from $I$ to $I$. A normal-play $P$-set admits no
$P$-to-$P$ edge, so $P\neq I$.
\end{proof}

\noindent\emph{Verification.} Cross-checked on the bent
$(8,1,\lambda=2)$ extraspecial group
(\path{experiments/gold/extraspecial_core.py}).

\subsection{Normal-play rigidity and the synthesis}

The strongest constraint is not any single architecture kill but a rigidity
statement inside a declared access class. Call a rule \emph{in-class} if the
quadratic refinement is quarantined behind a weight-bounded oracle (the
$(w_0,c)$ metering of Definition~\ref{def:criterion} below) and the rule is
\emph{refinement-uniform}: it realizes the target for every refinement $q$
of $B$ in the declared family, not just the Gold one.

\begin{theorem}[normal-play rigidity, in-class]\label{thm:rigidity}
Under the bounded-framing oracle model with refinement uniformity, every
legal edge between positions of Hamming weight $>c\,w_0$ must flip $Q$ --
in normal, misère, and loopy-Loss semantics alike. Consequently every
in-class realizer of $\Sigma$ in those semantics is a \emph{$Q$-alternating
ender}: play telescopes $Q$, and the outcome is the parity of $Q(x)$.
\end{theorem}

\begin{proof}[Proof sketch]
The case $Q(v)=Q(w)=0$ of a non-flipping edge is forbidden directly by the
outcome structure in all three readings. The case $Q(v)=Q(w)=1$ is reduced
to it by an adversary substitution: replace $q$ by $q'=q+\ell$ with $\ell$
a linear form annihilating every queried functional but not the dense
indicators; refinement uniformity transports the rule to $q'$, where the
edge becomes a forbidden $0$-to-$0$ edge.
\end{proof}

\begin{theorem}[no live middle, relative to the quarantine]\label{thm:nolivemiddle}
Let $Q$ have nonsingular core of rank $2r\geq4$. No fixed uniform
single-board rule satisfying
\begin{enumerate}[label=(\Alph*),leftmargin=*,itemsep=0pt]
\item[(S)] no affine digraph automorphism outside $\AO(Q)$,
\item[(A)] legality from $B$-oracle queries at $t\leq 2r-3$ frame constants,
\item[(C)] no commutative squaring route for the provenance of $Q$,
\item[(E)] no left-$E$-equivariant structure on the extraspecial group,
\item[(R)] bounded-framing oracle access with refinement uniformity,
\end{enumerate}
has $P$/Loss/Draw target equal to $\Sigma$ with any bulk strategic
content: under (R), every legal edge between dense positions flips $Q$, all
strategic content is confined to a Hamming ball of radius $c\,w_0$, and no
bulk-decision-live realizer exists.
\end{theorem}

\begin{remark}[scope of the no-go theorem]\label{rem:hatches}
Theorem~\ref{thm:nolivemiddle} is relative to its hypotheses. It does not
cover:
\begin{enumerate}[leftmargin=*,itemsep=0pt]
\item \emph{Loopy-Draw semantics}: the substitution contradiction of
Theorem~\ref{thm:rigidity} does not fire (Draw-to-Draw edges are legal).
\item \emph{$t\geq 2r-2$ with anisotropic complement}: escapes
Theorem~\ref{thm:F} at the symmetry level.
\item \emph{Frobenius-aware access}: the Galois group lies inside
$\Orth(Q_a)$, so all symmetry methods are silent;
Theorem~\ref{thm:F}'s model explicitly excludes Frobenius-twisted queries.
\item \emph{Non-quarantined rules using the game-native $\wp(w)$ diagonal
source} (Proposition~\ref{prop:wp}): once $q$ is game-built, refinement
uniformity holds trivially and class (R) says nothing.
\item \emph{Rank $r=1$ and radical-anisotropic degenerate layers}: the
character-sum count in Theorem~\ref{thm:H} vanishes, and small extraspecial
groups can be realized.
\end{enumerate}
The main construction uses the fourth route: a game-built diagonal source and
a public Witt-frame lift. Section~\ref{sec:criterion} replaces vocabulary
quarantine by the invariant transcript-span boundary, and
Section~\ref{sec:wittfifo} attains that boundary with an explicit local rule.
The list therefore records the theorem's scope, not unresolved premises of the
Gold realization.
\end{remark}

\section{Local realization and the observation boundary}\label{sec:criterion}

Two natural first attempts at ``naturality'' fail for structural reasons.
\emph{Equivariance fails}: by Theorem~\ref{thm:E} there is no invariance
class strictly between symmetry groups that kill the $P$-set and symmetry
groups that force extensional factoring through $Q$- and $B$-evaluation.
\emph{Vocabulary quarantine fails}: the Gold diagonal
$q_i^{(m,1)}=\Tr(e_i^3)$ is expressible as $\Tr(\wp(w)e_i)$ using only XOR,
nim-product, and trace (Proposition~\ref{prop:wp}) -- the licensed operation
chain itself -- so no principled line separates ``computing $q$'' from
``computing $Q$'' on vocabulary alone. The criterion below replaces both
with access metering and a sharp transcript-observation boundary.

\begin{definition}[uniform local realization]\label{def:criterion}
Fix the coin frame $\{e_i\}$ (the Grundy basis $g(n)=2^n$, game-natural).
\emph{Public data}: the frame, $\oplus$, nim-product, Frobenius, and the
full alternating form $B$ (game-built via Turning Corners, Frobenius, and
trace). \emph{Quarantined data}: the quadratic refinement, presented as a
weight-bounded oracle $O_q(z)=Q_q(z)$ for $\wt(z)\leq w_0$.
A \emph{uniform local rule} satisfies:
\begin{enumerate}[leftmargin=*,itemsep=0pt]
\item[\textbf{F1}] (\emph{$q$-blind statics}) position set, loading map,
and terminal set are computed without oracle access;
\item[\textbf{F2}] (\emph{metered access}) each move predicate and state
transition together use at most $c$ adaptive oracle queries, each at a point
of Hamming weight at most $w_0$, with $(w_0,c)$ constants independent of
$m$;
\item[\textbf{F3}] (\emph{declared semantics}) one canonical outcome
labeling (normal, misère, loopy, or $q$-blind terminal payoff).
\end{enumerate}
It is an \emph{exact local realizer} if additionally:
\begin{enumerate}[leftmargin=*,itemsep=0pt]
\item[\textbf{N1}] (\emph{torsor uniformity}) for every
$m$, every $B$ in the declared family, and every refinement $q$ of $B$:
$\mathrm{outcome}(\iota(x))=\text{target}$ iff $Q_q(x)=0$;
\item[\textbf{N2}] (\emph{locality budget}) the $(w_0,c)$ metering of F2;
\end{enumerate}
A rule with a reachable position having both a winning and a losing move is
called \emph{decision-nondegenerate}. These are objective properties of a
declared rule and oracle interface. They do not, by themselves, turn the
informal adjective ``natural'' into a semantic invariant.
\end{definition}

\begin{theorem}[Tier-3 exclusion under N1+N2]\label{thm:tier3excl}
For any position $x$ with $\wt(x)>c\,w_0$, the at most $c$ queried
functionals span only weight-$\leq c\,w_0$ directions, so there exist
refinements $q,q'$ agreeing with every oracle answer but differing on
$Q(x)$. Hence no in-class rule's move predicate factors through $Q(x)$ at
dense positions -- by theorem, not stipulation.
\end{theorem}

\begin{definition}[bounded outcome certificate]\label{def:certificate}
Write $\mathrm{N2cert}(C,w)$ when, for every refinement $q$ and input $x$,
there is a finite observation set $S(q,x)$ with $|S(q,x)|\leq C$ and
$\wt(z)\leq w$ for $z\in S(q,x)$ such that every same-polar refinement $q'$
agreeing with $q$ on $S(q,x)$ has the same rooted outcome. The set may be
chosen adaptively from $q$; what is bounded is the complete certificate for
the outcome, not the number of queries made by one transition.
\end{definition}

\begin{theorem}[sharp observation lower bound]\label{thm:certificate}
Let the refinement family at fixed $B$ be closed under addition of every
linear functional. If a transcript-stable rooted result is exact at $x$, then
\[
 x\in\operatorname{span} S(q,x).
\]
In the original frame, if every observed vector has weight at most $w$, then
\[
 \wt(x)\leq |S(q,x)|w.
\]
Consequently no constants $C,w$ give an exact torsor-uniform family satisfying
$\mathrm{N2cert}(C,w)$ in unbounded dimension.
\end{theorem}

\begin{proof}
If $x\notin\operatorname{span}S$, finite-dimensional separation gives a
linear functional $\ell$ vanishing on $S$ but not at $x$. The refinements
$q$ and $q+\ell$ therefore have identical observation transcripts and rooted
outcomes, while their target values at $x$ differ, contradicting exactness.
For the coordinate bound, the union of the supports of $|S|$ weight-$w$
vectors has size at most $|S|w$ and must cover the support of $x$.
Both forms are kernel-checked in
\path{formal/Ogdoad/GoldNoEvaluator.lean}.
\end{proof}

\begin{theorem}[fork padding obstruction]\label{thm:forkpadding}
The properties ``there exists a refinement-sensitive winning fork at a
reachable position'' and the stronger variant requiring that fork after every
original completed (hence every optimal completed) play do not certify that a
finite normal-play realizer is non-evaluative. Every finite normal-play tree
can be padded, without changing its root outcome, so that every completed play
enters an always-$N$ two-action fork whose unique winning action swaps with one
chosen refinement bit.
\end{theorem}

\begin{proof}
Replace every terminal $P$-node by a one-option wrapper whose child is the
always-$N$ swapping fork. The wrapper is still $P$, so induction over the
finite tree preserves every ancestor outcome. The new fork is forced after
every original completed play. The construction and induction are
kernel-checked as \texttt{win\_padTerminals} in
\path{formal/Ogdoad/GoldForkPadding.lean}.
\end{proof}

\begin{remark}[Why the old N3 screens are retired]\label{rem:N3}
The old mixed-successor N3 is defeated by dominated escape edges, while after
dominance pruning every retained successor at a position has one outcome
class. Counterfactual winning-action variants remain vulnerable to
Theorem~\ref{thm:forkpadding}, even if their witness is required on every
optimal play. Thus the paper uses the representation-invariant certificate
lower bound of Theorem~\ref{thm:certificate} for the negative boundary and
states the positive rule's syntax and access discipline directly. Strategic
liveness remains useful descriptive evidence, not an axiom defining
naturality.
\end{remark}

\section{The weighted-source Witt--FIFO construction}
\label{sec:wittfifo}

The arbitrary-graph linking theorem above is stronger than the Gold problem
requires. The polar form is public data, so it may be put into a deterministic
Witt frame before play. This changes the strategic support graph from an
arbitrary induced graph to a matching. A naive version would attach the dense
value $Q(f_i)$ to a Witt-basis coin $f_i$, violating N2 in the original
Grundy frame. The decisive point is to split that value into a public polar
part and weight-one source \emph{interactions}: the refinement bits weight
local overlap edges rather than being accumulated as invariant terminal
labels.

Fix the original frame $e_1,\dots,e_m$ and write
\[
 P_B(z)=\sum_{j<k}B(e_j,e_k)z_jz_k,
 \qquad Q_q(z)=P_B(z)+\sum_j q_jz_j,
 \qquad q_j=Q_q(e_j).
\]
Choose, by a fixed Gaussian-elimination convention depending only on $B$, a
symplectic-plus-radical basis
\[
 f_i=\sum_j C_{ji}e_j,
 \qquad B(f_{a_h},f_{b_h})=1
\]
whose other pairings vanish. For $x=\sum_i y_if_i$, put
\[
 p_i=P_B(f_i)=\sum_{j<k}C_{ji}C_{ki}B(e_j,e_k).
\]
All of $C$, $y=C^{-1}x$, and $p_i$ are refinement-blind public data.

For completeness, the adapted basis exists by the usual two-line induction for
alternating forms. If $B=0$, retain any basis as radical vectors. Otherwise
choose the lexicographically first pair $u,v$ with $B(u,v)=1$. The plane
$H=\langle u,v\rangle$ is nonsingular, since a vector $au+bv$ orthogonal to
both $u$ and $v$ has $a=b=0$; hence $V=H\oplus H^\perp$. Recurse on
$H^\perp$. Choosing the first available pivot and performing row operations in
the fixed original-coordinate order makes the construction deterministic. It
uses only the public matrix $B$, works unchanged in the singular case, and
terminates after one dimension drop for a radical step or two for a
hyperbolic-plane step.

\begin{quote}
\textbf{Rule (weighted-source \textsc{witt-fifo}).}
Load a strategic coin $F_i$ for each $y_i=1$, with fixed public diagonal
charge $p_i$. Join $F_{a_h}$ to $F_{b_h}$ exactly when both are loaded.
Every loaded strategic Witt edge has weight one.
For each original coordinate $x_j=1$, load in addition a q-blind source pair
$S_j^0,S_j^1$. Its potential edge has local weight
$q_j=O_q(e_j)$: opening one endpoint while its mate is already queued xors
$q_j$ into the charge. Source closes have charge zero.

Every loaded coin must be touched twice. An untouched coin may be opened;
only the FIFO front may close; opening onto an empty queue sets the one-move
ko; and a player passes exactly when no open or close is legal, clearing ko.
Initialize $\sigma=0$. When a coin $v$ opens, update $\sigma$ by the weight
of its already-open mate, if any. Thus every strategic or source edge is
charged exactly when its two touch intervals overlap. When a strategic front
$F_i$ closes, update
\[
 \sigma\mathrel{\oplus{=}}p_i.
\]
All other closes and all passes leave $\sigma$ unchanged.
After the queue is drained, add one ordinary normal-play claim move exactly
when $\sigma=1\oplus\phi$, where $\phi$ is the current move-parity bit
initialized to $0$ and toggled after every move.
\end{quote}

For a Gold form, Theorem~\ref{thm:all-diagonal-source} supplies
$q_j=\Tr((w_a^2+w_a)e_j)$ for the source-pair weight. Definition~\ref{def:criterion}
counts this as one weight-one oracle observation on an overlapping source
OPEN. It does not claim that constructing $w_a$ by exact finite-field linear
algebra, or evaluating its trace, is itself one primitive game move.

The move-parity bit is not an evaluator or form oracle; it is already the
mover field of the FIFO state. It is necessary for an identity-valued payoff
compiler: a phase-free terminal leaf assigns opposite initial-seat winners at
even and odd depths. The final claim is therefore a two-state local ``charge
lantern'', not an inspection of $x$, $B$, or $Q(x)$.

\begin{theorem}[FIFO matching theorem]\label{thm:matchingfifo}
Let every component of $G$ be $K_1$ or $K_2$. From the empty-queue root,
either designated seat has a strategy forcing the sum of all close charges
$\deg_U(f)$ to be zero. In fact the strategy makes every close charge zero
individually. The theorem holds for every finite board and needs no dummy.
\end{theorem}

\begin{proof}
Call a FIFO front \emph{safe} when either ko is set or its unique mate is not
untouched; equivalently, every legal close at that checkpoint has charge zero.
Induct on $4|U|+2|\mathcal Q|+[\mathrm{ko}]$. We show that a designated seat
forces terminal score zero whenever it moves from score zero, and whenever its
opponent receives a score-zero safe state.

If the designated seat moves with $U\ne\varnothing$ and the queue is empty,
OPEN any $z\in U$; ko makes the next checkpoint safe. If the front $f$ has
charge one, its unique mate $u$ lies in $U$; OPEN $u$, which erases the only
possible live neighbour of $f$. If $f$ has charge zero, OPEN any $z\in U$;
maximum degree one implies that erasure preserves the zero front. When
$U=\varnothing$, all later closes have charge zero and the forced PASS, when
present, is harmless.

At an opponent checkpoint, OPEN and PASS preserve the score and return the
turn. A legal CLOSE rules out ko, so safety makes its charge zero and likewise
returns a score-zero state. Every branch has smaller rank. Initially the queue
is empty: the second seat receives a vacuously safe root, while the first seat
uses the empty-queue OPEN. The designated seat itself never closes while an
untouched coin remains, and after the untouched set is empty every close has
charge zero. Thus the constructed strategy makes each close charge zero, not
only their total. This proves both seats.
\end{proof}

The exact transition-system proof is kernel-checked as
\texttt{evenWins\_initial\_of\_matching} in
\path{formal/Ogdoad/FifoMatching.lean}. The finite matching census helped find
the policy but is not used in the theorem. Its refinement-blind strengthening,
with a public potential matching and an arbitrary edge-deleted scoring
submatching, is kernel-checked as
\path{evenWins_initial_of_every_submatching}. The proposition is
pointwise in the scoring submatching, while its proof branches choose moves
only from the public matching. Since \texttt{EvenWins} does not reify a
first-class policy or complete play traces, the literal single-policy
quantifier and the per-close-zero observation are the explicit argument above,
not separate theorem statements in Lean.

\begin{theorem}[weighted-source exactness]\label{thm:wittfifo}
For every $\F_2$-valued quadratic refinement $Q_q$ on a finite-dimensional
$\F_2$-space, every input $x$, and either charge stance, weighted-source
\textsc{witt-fifo}
forces terminal charge $\sigma=Q_q(x)$. Consequently its charge-lantern
normal-play root is $P$ if and only if $Q_q(x)=0$.
\end{theorem}

\begin{proof}
Quadratic expansion in the adapted basis gives
\[
 P_B(x)=\sum_i y_iP_B(f_i)+\sum_{i<k}y_iy_kB(f_i,f_k)
 =\sum_i y_ip_i+\sum_h y_{a_h}y_{b_h}.
\]
The linear refinement part transports without any dense query:
\[
 \sum_i y_i\sum_j C_{ji}q_j=\sum_j x_jq_j.
\]
Therefore
\begin{equation}\label{eq:splitsource}
 Q_q(x)=\sum_i y_ip_i+\sum_jx_jq_j+\sum_hy_{a_h}y_{b_h}.
\end{equation}
Form instead the q-blind \emph{potential matching} $P$: retain every loaded
strategic Witt edge and every loaded source-pair edge, irrespective of its
weight $q_j$. Apply the explicit safe-front strategy of
Theorem~\ref{thm:matchingfifo} to $P$. The policy consults only this potential
matching and the public FIFO state: when the front's potential mate is still
untouched it opens that mate, and otherwise it opens any untouched coin. It is
therefore the same strategy for every refinement $q$. The proof of the theorem
shows more than cancellation of a total score: every close has zero live
$P$-degree.

For one potential edge $e$, exactly one of the following occurs: its two touch
intervals overlap, or the earlier-opened endpoint closes while the other
endpoint is still untouched. The latter case would give that close live
$P$-degree one, so the strategy excludes it. Consequently every potential
pair overlaps. The rule's total OPEN charge is therefore directly
\[
 \sum_h y_{a_h}y_{b_h}+\sum_jx_jq_j,
\]
with no refinement-dependent auxiliary graph or strategy.
Every active strategic $F_i$ closes exactly once and supplies
$\sum_i y_ip_i$. Equation~\eqref{eq:splitsource} therefore gives forced
terminal charge $Q_q(x)$ at both stances.

For the final assertion, retain the finite charge game tree. At a leaf of
charge $\sigma$ and phase $\phi$, the unique claim exists exactly when the
player to move is the seat designated by $\sigma$. Backward induction says
that the compiled normal-play winner equals the charge-game winner at every
subtree and phase. At the root, seat zero plays for charge one, so the root is
$N$ for $Q_q(x)=1$ and $P$ for $Q_q(x)=0$. The general compiler and root
equivalence are kernel-checked in \path{formal/Ogdoad/GoldSemantics.lean}.
\end{proof}

\begin{proposition}[locality and optimal observation support]
\label{prop:wittfifonatural}
The rule satisfies F1--F3 and N1--N2 with $(w_0,c)=(1,1)$. Its distinct
refinement observations at input $x$ are exactly the singleton directions
$\{e_j:x_j=1\}$; they span $x$ and attain the weight-one lower bound of
Theorem~\ref{thm:certificate}.
\end{proposition}

\begin{proof}
The basis algorithm, $x\mapsto y$, both loaded support sets, the matching
potential edges, public charges $p_i$, and terminal board condition use only
$B,x$ and the fixed frame. Thus F1 holds. An OPEN involving a source pair
needs at most the one singleton value $O_q(e_j)$, and only when its mate is
already queued; every other move predicate and transition is
refinement-blind. Even the forcing strategy is refinement-blind: it is the
safe-front policy on the public potential matching $P$. This proves F2 and N2
with $(1,1)$. The semantics is the single normal-play claim rule, giving F3, and
Theorem~\ref{thm:wittfifo} is N1 for every refinement of $B$.

There is one source pair exactly when $x_j=1$, so the set of possible oracle
directions is exactly $\{e_j:x_j=1\}$. These vectors are independent and sum
to $x$; no proper subset spans $x$. Theorem~\ref{thm:certificate} says every
transcript-stable exact weight-one interface must cover the same support.
\end{proof}

The oracle notation is accounting, not hidden information: $q$ is a public
parameter of the perfect-information game. The play-dependent source-pair
repair is essential.  A rule using only invariant source charges would make
every complete history have the form
\[
 \sigma_q(h)=c_{B,x}(h)+\sum_{j:x_j=1}q_j,
\]
so two refinements agreeing on $Q(x)$ produced action-labelled isomorphic game
trees. In the present rule the coefficient of $q_j$ is the play-dependent
overlap of $S_j^0,S_j^1$; refinement data therefore enters the same local
interaction law as the polar edges, not a history-independent terminal bit.
This is an intensional design fact, not an outcome-only certification:
Theorem~\ref{thm:forkpadding} proves that the reachable and unavoidable fork
properties do not provide such a certification.

The rank-zero boundary is left undecorated. The same weighted-source rule still
realizes the linear refinement exactly; there simply are no strategic Witt
edges to pretend otherwise.

The original isolated-dummy arbitrary-graph linking conjecture remains an
interesting strict generalization, but it is no longer load-bearing for Gold:
the public Witt frame reduces precisely the required boards to the proved
matching class, while weighted source pairs preserve fixed-frame locality.

\section{Formal boundary and extensions}\label{sec:status}

The construction separates into six independently checked components:
\begin{itemize}[leftmargin=*]
\item \path{formal/Ogdoad/GoldDiagonal.lean}: trace blocks and the
      Artin--Schreier identities behind the all-exponent source;
\item \path{FifoMatching.lean}: the both-seat matching strategy;
\item \path{GoldMatchingAlgebra.lean}: the adapted-basis expansion;
\item \path{GoldSemantics.lean}: the phase-aware claim compiler;
\item \path{GoldNoEvaluator.lean}: the observation lower bound;
\item \path{GoldForkPadding.lean}: outcome-preserving padding.
\end{itemize}
The paper supplies the concrete Witt basis, the original-frame coefficient
transport, and the synthesis of those components. There is no single
end-to-end Lean theorem constructing the complete arena.

The arbitrary-graph isolated-dummy FIFO theorem remains open and is studied in
\path{writeups/linking_affine.tex}. It would strengthen FIFO combinatorics but
does not strengthen Theorem~\ref{thm:wittfifo}, whose boards are matchings plus
isolates. The all-exponent diagonal source is supplied by
Theorem~\ref{thm:all-diagonal-source}; it is not an additional hypothesis.

\section{Conclusion}

Gold trace forms therefore have a uniform normal-play semantics built from
nimber operations and local interactions. A public Witt frame makes the polar
graph a matching, weighted source pairs preserve original-frame singleton
locality, the safe-front theorem forces the exact charge in every dimension,
and the phase-aware claim compiles that charge to a $P/N$ outcome. The Arf
invariant and radical data give the resulting zero-set bias.

The surrounding negative results make the access boundary sharp. Linear
coin-turning reaches the lexicode floor but not a genuine quadratic zero set;
the extraspecial extension identifies the missing noncommutative datum;
symmetric, commutative, and bounded-total-observation architectures cannot
carry it uniformly; and extensional fork tests cannot define naturality.
Weighted-source Witt--FIFO meets the necessary observation span with singleton
support. The remaining arbitrary-graph FIFO contraction is a separate
combinatorial conjecture, not a gap in the Gold theorem.

\begin{thebibliography}{9}
\bibitem{Arf41} C.~Arf, \emph{Untersuchungen \"uber quadratische Formen in
K\"orpern der Charakteristik~2, I}, J.\ Reine Angew.\ Math.\ \textbf{183}
(1941), 148--167.
\bibitem{Bouton} C.~L.~Bouton, \emph{Nim, a game with a complete mathematical
theory}, Ann.\ of Math.\ \textbf{3} (1901/02), 35--39.
\bibitem{WW} E.~R.~Berlekamp, J.~H.~Conway, R.~K.~Guy, \emph{Winning Ways for
Your Mathematical Plays}, 2nd ed., A~K Peters, 2001--2004.
\bibitem{Conway76} J.~H.~Conway, \emph{On Numbers and Games}, Academic Press,
1976.
\bibitem{ConwaySloane86} J.~H.~Conway, N.~J.~A.~Sloane, \emph{Lexicographic
codes: error-correcting codes from game theory}, IEEE Trans.\ Inform.\
Theory \textbf{32} (1986), 337--348.
\bibitem{Dickson01} L.~E.~Dickson, \emph{Linear Groups, with an Exposition of
the Galois Field Theory}, Teubner, 1901.
\bibitem{EKM08} R.~Elman, N.~Karpenko, A.~Merkurjev, \emph{The Algebraic and
Geometric Theory of Quadratic Forms}, AMS Colloquium Publications
\textbf{56}, American Mathematical Society, 2008.
\bibitem{Gold68} R.~Gold, \emph{Maximal recursive sequences with $3$-valued
recursive cross-correlation functions}, IEEE Trans.\ Inform.\ Theory
\textbf{14} (1968), 154--156.
\bibitem{Gorenstein80} D.~Gorenstein, \emph{Finite Groups}, 2nd ed.,
Chelsea Publishing, New York, 1980.
\bibitem{LN} R.~Lidl, H.~Niederreiter, \emph{Finite Fields}, 2nd ed.,
Cambridge University Press, 1997.
\bibitem{Ovsienko16} V.~Ovsienko, \emph{Real Clifford algebras and quadratic
forms over $\F_2$: two old problems become one}, The Mathematical Intelligencer
\textbf{38} (2016); arXiv:1601.07664.
\bibitem{PS} T.~E.~Plambeck, A.~N.~Siegel, \emph{Misere quotients for impartial
games}, J.\ Combin.\ Theory Ser.\ A \textbf{115} (2008), 593--622.
\bibitem{Quillen71} D.~Quillen, \emph{The mod~$2$ cohomology rings of
extra-special $2$-groups and the spectra of quadratic forms}, Math.\ Ann.\
\textbf{194} (1971), 197--212.
\bibitem{Taylor} D.~E.~Taylor, \emph{The Geometry of the Classical Groups},
Sigma Series in Pure Mathematics~\textbf{9}, Heldermann Verlag, 1992.
\bibitem{Winter72} D.~L.~Winter, \emph{The automorphism group of an
extraspecial $p$-group}, Rocky Mountain J.\ Math.\ \textbf{2} (1972),
159--168.
\end{thebibliography}

\end{document}
